\section{Methods and data}\label{sec:methods}

\subsection{Metric}
We compare options by the cost of the electricity that actually meets demand in every period, rather than by the levelised cost of the energy each plant produces. For a single technology on a monthly basis this is the levelised cost multiplied by the overbuild needed to cover the worst month, plus the cost of the storage that shifts energy into the hours and months in which it is missing (Section~\ref{sec:monthly}). For the system as a whole it is the total annualised cost of the least-cost portfolio that meets demand in every time step, divided by the demand not covered by hydropower, \texttt{cost\_new} in the repository. This single denominator (356.6\,TWh of the 400\,TWh) is used throughout the system results; multiplied by 356.6/400 = 0.8915 it gives the cost per MWh of total demand, for example <<EXPR:V('RUN:ss0_PI_670_h3:cost_new')*0.8915|1>> instead of <<RUN:ss0_PI_670_h3:cost_new|2>>\,EUR/MWh for the central PI case at three hours. The monthly calculation divides instead by total demand. Monthly and hourly costs must be compared only through this conversion. Because the time step biases absolute costs (Section~\ref{sec:step}), the decision metric is the gap, the cost of the system without nuclear minus the cost of the system with nuclear.

\subsection{Capacity expansion model}\label{sec:model}
The model is a linear programme that chooses the capacity of nuclear ($N$), of each block $b$ of the solar ($S_b$), onshore wind ($W_b$) and offshore wind ($O_b$) supply curves, and the energy capacity of lithium-ion batteries ($E$), and their dispatch, so as to meet demand in every time step at least cost, with zero direct emissions (mode ZD). Version 7 adds every other technology that can produce, store, shift or reduce electricity in Italy as a decision variable with one uniform cost rule: imports (11.5\,GW existing and 7.3\,GW new), existing and new hydropower and pumped hydro, geothermal (existing, two conventional blocks and enhanced systems), bioenergy (biogas, solid biomass, waste), concentrated solar power with storage, wave energy, compressed air, hydrogen (electrolysis, geological storage, turbine), 100-hour storage, smart charging, interruptible load and small modular reactors; gas with carbon capture and direct air capture enter only in the net-zero (NZ) variants. No technology is excluded without a number showing that it is negligible (tidal currents, for example, are 0.031\% of demand). Technologies with market prices or contracts use the latest observed price. Technologies without commercial plants have two cases, both reported: PI (first plant, costs of demonstrators or first offers) and OB (target, costs declared by developers or projected). The OB model is reduced: it has no wave energy, 100-hour storage, new pumped hydro, new hydropower or SMR. The time step is six hours for the curve points and three hours for the corrected central values (Section~\ref{sec:step}); the 3-hour data are block sums, so a one-hour step is not possible. The four weather years 2016--2019 are optimised jointly: one set of capacities must serve all of them, with dispatch optimised in each year. A power system is built before the weather is known, and optimising each year on its own would give every year a fleet designed for its own weather. The model has a single node, so energy moves between regions without losses or constraints beyond a grid charge on capacity, and it has no ramp constraints: nuclear output may move between its minimum and its maximum from one step to the next at no cost. The objective is
\begin{equation}
\min \; Y\Big(c_N N + \sum_b c_{S_b} S_b + \sum_b c_{W_b} W_b + \sum_b c_{O_b} O_b + c_E E\Big) + \sum_{t} v_N\, n_t ,
\end{equation}
where $Y$ is the number of years optimised together and $c_i$ the annualised cost of capacity $i$, including the grid charge. For solar and wind the cost of a block is the real lifetime cost derived from auction prices (Section~\ref{sec:inputs}) multiplied by the full-load hours of the block, so that each MWh the block can produce is paid whether or not it is used. Existing solar plants installed after 2016 and the existing wind block E0 enter at operating cost only. The energy balance in each step $t$ of length $\Delta$ is
\begin{equation}
n_t + p^S_t + p^W_t + p^O_t + h_t + d^B_t - q^B_t + d^P_t - q^P_t = L_t ,
\end{equation}
where all flows are energies in step $t$, $L_t$ is the demand net of the flat generation from bioenergy and geothermal, $n_t$ nuclear output, $p^S_t \le \sum_b S_b\,\rho_b\,s_t$ the solar output used ($s_t$ the PECD energy per MW and $\rho_b$ the scale of the block for its zone), with the same form for onshore and offshore wind, so that curtailment is allowed; $h_t$ is hydropower and $d$, $q$ are discharge and charge of batteries ($B$) and pumped hydro ($P$). Nuclear output follows load between a minimum and a maximum share of the available capacity,
\begin{equation}
\mu\,a_t\,N\Delta \le n_t \le a_t\,N\Delta ,
\end{equation}
with minimum load $\mu=0.5$ and availability $a_t$ that depends on the month: 0.97 in December, January, February, June, July and August and 0.87 in the other six months, a mean of 0.92. This availability is an assumption: it is built to match the 92\% load factor of the US fleet \citep{EIA2025} and refuelling outages concentrated in spring and autumn \citep{EIA2023}, not a measurement of an Italian fleet. Storage follows
\begin{equation}
e^B_{t+1} = e^B_t + \eta_B q^B_t - d^B_t,\quad 0\le e^B_t \le E,\quad q^B_t, d^B_t \le E\Delta/\tau ,
\end{equation}
with round-trip efficiency $\eta_B=0.85$ and duration $\tau=6.6$\,h, and the same form for pumped hydro with fixed power 7.3\,GW, energy 53\,GWh and $\eta_P=0.75$; the other storage technologies and imports enter the balance with their own flows and costs (Table~\ref{tab:inputs}). Storage levels are cyclic. Hydropower in each month $m$ must produce the monthly energy $H_m$, with a minimum equal to a share $f=0.37$ of the monthly mean (run-of-river flow \citep{TernaHydro}) and a maximum power $\bar h=15$\,GW; the budgets $H_m$ are the same in the four simulated years, so hydropower is independent of the simulated wind and solar weather:
\begin{equation}
\sum_{t\in m} h_t = H_m,\qquad f\,\frac{H_m}{|m|}\le h_t \le \bar{h}\Delta .
\end{equation}

\paragraph{Adequacy reserve.} A static constraint requires the credited capacity to exceed the net peak $\hat L$ (66.98\,GW at six hours, 76.88\,GW at three hours) by a factor $\rho=1.1$:
\begin{equation}
a_{\mathrm{pk}} N + E/\tau + \bar h_c + \bar P + \kappa I \ge \rho \hat L ,
\end{equation}
where $a_{\mathrm{pk}}=0.97$ is the nuclear availability in the month of the peak, $\bar h_c=15$\,GW the hydro credit, $\bar P=7.3$\,GW the pumped hydro, $I$ the import capacity and $\kappa$ its credit: zero in the central case, 0.38 in the variant of the earlier version (4.4 of 11.5\,GW in the 2025 capacity market); credits of the other technologies are in the repository. All credits are declared assumptions.

\paragraph{Imports and net self-sufficiency.} Imports $i_t$ use the interconnection capacity $I$ (11.5\,GW existing and up to 7.3\,GW new) at 115.94\,EUR/MWh, the 2025 average PUN \citep{GME2025}. Exports $x_t\le I\Delta$ use the same capacity and earn nothing, a cautious assumption for the system that exports, because surplus energy is worth little precisely when it is also abundant in neighbouring countries. The central rule is net self-sufficiency: for every weather year $y$,
\begin{equation}
\sum_{t\in y} i_t \le \sum_{t\in y} x_t + X\,D_y ,\qquad X=0 ,
\end{equation}
where $D_y=400$\,TWh, and in the three critical weeks of each year imports are not available and $\kappa=0$. The rule asks that the Italian plan would still work if every European country planned the same way: an Italy that covers a fixed share of demand with energy produced abroad leans on plants that others build and pay for, much of them French nuclear. Variants: $X=5\%$ and $10\%$; imports without an annual bound and without exports but with no help in the critical weeks; and the rule of the earlier version, imports without exports, a 17\% price mark-up in the critical weeks and an availability and credit of 38\%. With $X=100\%$ and the 38\% credit the model reproduces the earlier central case to the cent. The model does not simulate the neighbouring systems. The twenty highest net-load steps all fall in winter evenings between 17:00 and 21:00, mostly in January, which is why heat pumps and electric vehicles push towards firm capacity.

\paragraph{Nuclear cost.} The annualised fixed cost combines overnight construction cost $K$, real cost of capital $r$ and construction time $T_c$:
\begin{equation}
c_N = K\,\phi(r,T_c)\,\mathrm{CRF}(r,L_N) + F_N ,\qquad
\phi(r,T_c)=\frac{1}{T_c}\sum_{k=0}^{T_c-1}(1+r)^{T_c-k-1/2},\qquad \mathrm{CRF}(r,n)=\frac{r(1+r)^n}{(1+r)^n-1},
\end{equation}
where $\phi$ accounts for interest during construction with uniform spending, $L_N=80$ years, and $F_N$ is the fixed part: operation at 17\,EUR/MWh at 90\% load, $17\times0.9\times8760=134.03$\,EUR/kW-yr, plus a grid charge of 2\,EUR/MWh, 16.11\,EUR/kW-yr, so $F_N=150.14$ at 5\% and 154.77 at 7\%, because at 7\% the grid charge is multiplied by 1.287. We call $c_N$ \emph{fixN}. The inverse is $K=(\mathrm{fixN}-F_N)/(\phi\,\mathrm{CRF})$. The variable cost is 15\,EUR/MWh. At 7\% the capital part of every technology priced by auction or by capex is multiplied by $\mathrm{CRF}(7\%,L)/\mathrm{CRF}(5\%,L)$ for its life $L$ (1.176 for 20 years, 1.239 for 30, 1.287 for 40); operation, fuel and the price of imported energy are unchanged, and nuclear uses its own $r$. The model is solved with HiGHS through its Python interface; all solutions have status optimal. The validation set is five reference runs (no-nuclear PI and OB, PI and OB at fixN 598 with imports, PI at 817 with 7\%), reproduced within 0.05\,EUR/MWh and 0.2\,GW.

\subsection{Data and assumptions}\label{sec:inputs}
Table~\ref{tab:inputs} lists the inputs with their status: DATA (read in a source), DERIVED (calculation on data), ESTIMATE (the author's judgement with a declared method) or ASSUMPTION (a model choice without a source, with a sensitivity).

\paragraph{Demand.} The 2040 demand is 400\,TWh without electrolysers. It is the sum of the ENTSO-E hourly load of 2016--2019 \citep{OPSD2020}, rescaled to 338\,TWh in each year, 15\,TWh of heat pumps and 47\,TWh of electric vehicles. The value is a central choice between the 385\,TWh for 2035 and the 415 and 439\,TWh for 2040 of the Terna and Snam scenario report \citep{DDS2024}. Heat pumps follow Italian monthly heating degree-days (Eurostat, \texttt{nrg\_chdd\_m}); their hourly shape, with morning and evening peaks, is an assumption. Electric vehicles charge 70\% of their daily energy between 18:00 and 24:00 and 30\% spread over the day, an assumption tested with flat, night-only, evening-only, daylight (50\% in daylight hours) and mixed profiles. Bioenergy (8\,TWh net) and geothermal (5.3\,TWh net) are subtracted as flat generation.

\paragraph{Weather.} Solar, onshore and offshore wind capacity factors are hourly PECD 2021.3 profiles for the 2030 fleet in Italy \citep{PECD2021}, used without normalisation: the mean full-load hours of 2016--2019 are 1{,}375 for solar, 2{,}227 for onshore wind and 2{,}935 for offshore wind, and differences between years are kept. For the worst-weather check we use the 38 PECD years 1982--2019 (Section~\ref{sec:results}).

\paragraph{Hydropower.} Annual hydro energy is 43.4\,TWh, the 2016--2019 mean excluding pumping, with the monthly shape of the 2023--2024 mean \citep{Terna2025m}. Pumped hydro power is 7.3\,GW (Terna 2024) and its energy 53\,GWh (Ambrosetti-Edison); efficiency is an assumption.

\paragraph{Solar supply curve.} Solar has 13 blocks by type (ground-mounted, industrial rooftops, other rooftops, agrivoltaic) and zone (South and islands, Centre-South, North), with a limit in GW, a scale $\rho_b$ for the yield of the zone from PVGIS 5.3 \citep{PVGIS}, and a cost multiplier on the capital part. Ground-mounted modules are banned on agricultural land by Article 5 of Decree-Law 63/2024 \citep{DL63}, a rule upheld by Constitutional Court judgment 127/2026 and taken up by Article 11-bis introduced by Decree-Law 175/2025 \citep{Consulta127}. No national estimate of the remaining non-agricultural area exists, so the limits are our estimates anchored to ISPRA \citep{ISPRA2026}: in the central curve 42.7\,GW on the ground, 140\,GW of rooftops (the lower end of the 140--193\,GW range of ISPRA 2026, of which 30\,GW industrial at 1{,}000\,EUR/kW) and 119.7\,GW of agrivoltaic (3\% of suitable agricultural area at 0.5\,MW per hectare), 302\,GW in all. For comparison \citet{Ranjgar2025} estimate 28--52\,GW of ground-mounted potential with 1\% of the available area. Ground-mounted plants occupy 1{,}700\,ha/GW \citep{ISPRA2026}. Output per kW for new plants is 10--16\% higher in the South than in the North (PVGIS, Sicily to Lombardy 1.16, Apulia to Veneto 1.10); the larger difference in the yield of the existing fleet mixes orientation, age and curtailment \citep{GSE2025}. The installed fleet was 43.5\,GW at the end of 2025 and 19.3\,GW at the end of 2016 \citep{ItaliaSolareQ4}; the 24.2\,GW installed after 2016 enter at operating cost and are subtracted from the blocks. The capital share of the real cost of solar is 76.4\%; multipliers apply to that part: 1.333 for agrivoltaic (1.0 against 0.75 million EUR/MW, midpoint of the GSE range reported by \citet{QualEnergiaAgri}) and 1.5625 for rooftops other than industrial, within the 900--1{,}600\,EUR/kWp range of \citet{FraunhoferEN2024}.

\paragraph{Wind supply curve.} Onshore wind has six blocks: E0, 13.6\,GW of existing repowered plants at operating cost (31\,EUR/MWh, 1{,}600 h), E1 6\,GW at 64\,EUR/MWh (2{,}400 h), E2 10\,GW at 71 (2{,}100 h), E3 15\,GW at 82 (1{,}800 h), E4 15\,GW at 96 (1{,}500 h) and E5 10\,GW at 118 (1{,}200 h), 69.6\,GW in all; block scale $\rho_b$ is hours of the block over the 2{,}227 h of the PECD profile. Costs follow JRC-EU-TIMES parameters (1{,}320\,kEUR/MW, 5.2\%, 25 years) \citep{JRCEUTIMES}, checked against the FER~X auction at 72.85\,EUR/MWh \citep{FERXmain2025}. The potential for Italy in the same source ranges from 10.1 to 171.6\,GW with the planning rules used. Offshore wind has five blocks, O1 2\,GW at 160\,EUR/MWh, O2 3.8\,GW at 172.5, O3 6\,GW at 138.5, O4 10\,GW at 112.5 and O5 25\,GW at 132.5, 46.8\,GW in all, with grid connection and export cable included in the cost and so no base grid charge.

\paragraph{Renewable and storage prices.} The first FER~X auction (December 2025) awarded 7.7\,GW of solar at an average of 56.82\,EUR/MWh \citep{FERXmain2025}. For batteries we use the first MACSE auction of 2025, 12{,}959\,EUR per MWh of capacity per year for 6.6-hour lithium-ion batteries and 15-year contracts \citep{MACSE2025}. The MACSE premium is the fixed remuneration paid to the holder of the capacity: market revenue from time shifting goes to the operator Terna and reduces the burden on consumers, and the holder keeps only a share of the margins from ancillary services \citep{ARERA247}; the premium is therefore a good estimate of the full capacity cost, if anything slightly low. FER~X prices are indexed to inflation for 20\% of the strike price after commissioning (art.\ 11 c.\ 3 of DM 457/2024) \citep{FERXrules2025,FERXindex}; the 20\% applied to every technology is an assumption, and the MACSE premium is indexed for 10\% \citep{MACSErules}. We convert each auction price $p$ into a real levelised cost over the plant life,
\begin{equation}
c = p\,\frac{\sum_{t=1}^{T}\left[\iota+(1-\iota)(1+\pi)^{-t}\right](1+r)^{-t}}{\sum_{t=1}^{T}(1+r)^{-t}}\left[(1-o)\frac{\mathrm{CRF}(r,L)}{\mathrm{CRF}(r,T)}+o\right],\label{eq:conv}
\end{equation}
with contract length $T$, indexed share $\iota$, inflation $\pi=2\%$, real discount rate $r=5\%$ for renewables and storage (also in the cases where nuclear is financed at 4\%), operating-cost share $o$ and plant life $L$. The real cost of solar is 42.2\,EUR/MWh, composed of capital (76.4\%) and operation (9.94\,EUR/MWh), and of batteries 11{,}466\,EUR/MWh-yr; indexation of 0\% or 100\% would give 40.7 or 48.2\,EUR/MWh (our calculation). Disposal is assumed to be in the offer price, because the rules put it on the producer or operator: the electronic-waste rules for solar \citep{RAEE}, the restoration guarantee required by the permit for wind \citep{WindRestoration} and extended producer responsibility for batteries \citep{BatteryReg}; the estimated effect is 1--2\,EUR/MWh. Plants are rebuilt every 30 years (solar and wind) and 15 years (batteries). Existing plants after 2016 and E0 are assumed not to need renewal before 2040.

\paragraph{Grid.} The grid charge is on new capacity: 13\,EUR/MWh for solar and 8 for onshore wind (a charge per MW-year equal to those values times 1{,}066 and 1{,}709 h) and 2\,EUR/MWh for nuclear, doubling on new solar and wind above 137\,GW, the capacity enabled by the 2035 transmission plan. The values derive from Terna's development plan, 23 billion euros for 94\,GW of new capacity, about 245\,EUR/kW, with a 5\% annuity over 40 years, our rounding of the ARERA conventional lives of 45 years for transmission lines and 33 for stations \citep{ARERA615}, applied to the plan of \citet{TernaPdS2025}; the doubling follows the increasing integration cost of \citet{Ueckerdt2013}, and the factor 2 is an estimate. At 7\% the grid charge is multiplied by 1.287 (our assumption for the scale of the grid).

\paragraph{Nuclear: cost cases.} Cost cases are points on one curve of fixN, not separate runs: Saeul 2{,}965\,EUR/kW with 8 years (fixN 335.2 at 5\%), Barakah 4{,}670 with 8 years (<<FIX:4670:5:8|1>>), the central case Dukovany 8{,}540 with 7 years (669.5 at 5\%, 922.7 at 7\%; the runs are at 670 and 923), a risk case 9{,}680 with 15 years (878.3 at 5\%, 1{,}334.2 at 7\%) and Vogtle 15{,}540 with 10 years (1{,}172.2). The central case stands for plants built properly, with one series design and engineering completed before construction starts; Flamanville, Olkiluoto 3, Mochovce and Vogtle are the Western first-of-a-kind builds behind the risk case, and Hinkley Point C is another risk reference that is not in its median. All five cases are overnight construction costs, without interest during construction (which the model adds through the factor idc), in 2025 euros per kW and including owner's costs. Dollar values are brought to 2025 with the US consumer price index (CPI-U, 274.310 in September 2021 and 321.943 on average in 2025, eleven months, October is missing because of the federal funding lapse \citep{BLSCPI}) and converted at 1 EUR = 1.1188 USD, the Bank of Italy reference rate of 26 October 2026 \citep{BancaItaliaFX}. The annual CPI-U averages before 2016, needed for Barakah and Vogtle, are the means of the twelve monthly values of the not seasonally adjusted series (BLS CUUR0000SA0, retrieved from FRED \citep{FREDCPI}): 214.537 in 2009, 218.056 in 2010, 224.939 in 2011, 229.594 in 2012, 232.957 in 2013, 236.736 in 2014 and 237.017 in 2015. The mean factors for Barakah and Vogtle recomputed with them are 1.2549 and 1.2623, equal to the values used (1.255 and 1.262), so no value changes.

\paragraph{Nuclear: Saeul and Barakah.} Saeul is the Korean APR1400 series (about 3{,}300\,USD/kW for a 1.4\,GW unit \citep{HSF2026,KHNP2025}). Barakah is the same design exported to the United Arab Emirates: the construction contract is about 20.4 billion USD for four units, 5.6\,GW \citep{WNAuae}, that is 3{,}643\,USD/kW nominal. The 24.4 billion USD of total financing in the same source includes interest and non-construction items and is not used, because the model adds interest itself. With uniform spending over 2012--2024 (mean CPI-U factor 1.255) the contract is 4{,}571\,USD/kW in 2025 dollars, 4{,}086\,EUR/kW. As at Dukovany the contract is a construction price (it also includes the first fuel loads, which we do not remove) and does not include owner's costs; adding them with the rule below gives 4{,}673, stated as 4{,}670\,EUR/kW (our calculation), 7\% above the 4{,}371 that the total financing would give.

\paragraph{Nuclear: the central case, Dukovany.} The central value is built in four steps. First, the contract price: the Czech Republic signed with KHNP on 4 June 2025 the contract for two APR1000 units of 1{,}055\,MWe \citep{Dukovany2024,WNNDukovanySigned2025,EnerdataDukovany2025}, CZK\,407 billion, USD\,18.6 billion, for construction, commissioning and fuel for about ten years (first load and five reloads) according to World Nuclear News \citep{WNNKHNPplans}. The exchange rate is the 2025 annual average of the ECB reference rate, 24.688 CZK per EUR \citep{EurostatFX}; an earlier version assumed 25, and the correction raises the cost by about 1.3\%. The price is then 16.5 billion euros, 7{,}813\,EUR/kW for 2 $\times$ 1{,}055\,MWe and 7{,}850 for 2{,}100\,MW; the model uses 7{,}850. The price is a June 2025 price, so it is in 2025 prices and is not escalated; the indexation in the contract concerns only the nominal price. Second, the fuel. The sources disagree. World Nuclear News says the CZK\,407 billion include fuel for about ten years (first load and five reloads) \citep{WNNKHNPplans}, while \v{C}T24, quoting the director of Elektr\'arna Dukovany II, and Forbes \v{C}esko report CZK\,15 billion more for the first fuel \citep{CT24Dukovany,ForbesDukovany}. We adopt the first because it states what the price covers (first load and five reloads), whereas the other two do not say whether the CZK\,15 billion are inside or outside the CZK\,407 billion. The central case therefore removes the fuel, because the model already counts fuel in the 32\,EUR/MWh of operating costs: fuel costs 4.85\,EUR/MWh (uranium, conversion, enrichment and fabrication), times 10 years, 0.9 load and 8{,}760 hours divided by 1{,}000, that is 382\,EUR/kW, and 7{,}850 minus 382 is 7{,}468 (our calculation). The fuel cost is 4.62\,USD/MWh in September 2021 prices \citep{WNAecon}; 4.62 times the CPI-U factor 1.174 divided by 1.1188 is 4.85\,EUR/MWh. The runs use 4.91, a negligible difference (0.4\,EUR/kW-yr in fixN). Effect of the other reading: if the CZK\,407 billion did not include the fuel and it were not removed, the cost before owner's costs would be 7{,}850 instead of 7{,}468 (382\,EUR/kW more) and the overnight cost 8{,}980 instead of 8{,}540 (440\,EUR/kW, 5\% more; fixN 696 instead of 670).

Third, owner's costs. Neither the contract nor the sources that report it \citep{WNNKHNPplans,CT24Dukovany,ForbesDukovany} give them, and we found no public EDU~II estimate that separates them, so the contract price is an EPC price and the overnight cost is higher. We add them with the share that the Idaho National Laboratory reports for a pressurised-water reactor in the EEDB cost database: owner's costs (account 40) are 12.57\% of overnight cost \citep{INLowner}, so overnight cost is the EPC price divided by $(1-0.1257)$, that is multiplied by 1.1438, which is 14.4\% of the EPC price. As a European reference, the synthesis by D'haeseleer for the European Commission (DG Energy, 2013) notes that there is no unique definition and reports from several sources owner's costs of about 15--20\% of the EPC cost \citep{DHaeseleer2013}: our 14.4\% lies just below the low end of that range, and its high end of 20\% gives 8{,}962\,EUR/kW (7{,}468 $\times$ 1.20), fixN 695. The ends of the range are read off the curve and are not rerun; the central value is the INL share. Fourth, the result: $7{,}468\times1.1438=8{,}542$, stated as 8{,}540\,EUR/kW (our calculation), with fixN 669.5 at 5\% and 922.7 at 7\%.

The central value has three bounds, all on the same curve. The lowest cost is the EPC price net of fuel alone, 7{,}470\,EUR/kW, fixN 604: optimal nuclear is about <<N6:PI:s:5:604|1>>\,GW (PI, six hours) against <<RUN:ss0_PI_670_h6:N_GW|1>> in the central case at the same step; it is the upper bound for nuclear because it ignores costs that exist. The highest cost adds the state preparation spending of CZK\,85 billion quoted by \v{C}T24 \citep{CT24Dukovany}, 21\% of the CZK\,407 billion, instead of the 14.4\% mark-up: 9{,}030\,EUR/kW, fixN 699 (close to the 20\% of D'haeseleer), with <<N6:PI:s:5:699|1>>\,GW in PI and <<N6:OB:s:5:699|1>> in OB at six hours, against <<N6:PI:s:5:670|1>> and <<N6:OB:s:5:670|1>> in the central case at six hours. This figure is an assumption: we do not know how much of the CZK\,85 billion is construction. Without removing the fuel (8{,}980\,EUR/kW, fixN 696) the six-hour result is <<N6:PI:s:5:696|1>>\,GW in PI and <<N6:OB:s:5:696|1>> in OB. Owner's costs are the least documented item of the paper, and the choice of the central case depends on them.

\paragraph{Nuclear: risk case and Vogtle.} The risk case, 9{,}680\,EUR/kW with 15 years, is the median of four recent Western builds in 2025 euros per kW: Mochovce 3--4 about 8{,}760, Olkiluoto 3 about 9{,}060, Flamanville 3 about 10{,}300 and Vogtle 3--4 15{,}540; the median is the mean of the two central values, (9{,}060 + 10{,}300)/2 = 9{,}680 (our calculation). Olkiluoto 3 cost about 11 billion \citep{NEIOL3} for 1{,}600\,MWe (1{,}575 net \citep{WNAfinland}), 6{,}875\,EUR/kW nominal. The source reads ``\$11bn (\$12bn)'': we read 11 billion euros, 12 billion dollars; if it were 11 billion dollars, Olkiluoto 3 would be about 8{,}100\,EUR/kW and the median about 9{,}530 (1.6\% lower). Mochovce 3--4 cost 6.7 billion euros \citep{WNNMochovce4} for 2 $\times$ 471\,MWe, 7{,}113\,EUR/kW nominal (7{,}614 with the 440\,MWe net per unit of \citealt{WNAslovakia}). Brought to 2025 euros with the HICP \citep{EurostatHICP}, assuming uniform spending in 2005--2022 for Olkiluoto and 2009--2025 for Mochovce (mean factors 1.318 and 1.231), they are about 9{,}060 and 8{,}760\,EUR/kW. The sources give neither the spending schedule nor the price year: with other schedules Olkiluoto ranges from 8{,}907 to 9{,}208\,EUR/kW. Flamanville cost 13.2 billion euros of 2015, excluding interim interest \citep{EDFFlamanville2022}, for 1.65\,GW gross \citep{Folz2019}: 8{,}000\,EUR/kW in 2015 euros, and with the euro-area HICP (2025 over 2015: 128.75 over 100.00, factor 1.2875 \citep{EurostatHICP}) about 10{,}300 in 2025 euros; it took about 17 years \citep{Flamanville2024}, and its cost per kW is about 21\% above the central case (10{,}300 over 8{,}540, our calculation). Vogtle 3--4 cost close to 35 billion USD in total \citep{EnergytechVogtle}, of which more than 3.5--4 billion are financing charges \citep{IEEFAVogtle}, so about 31 billion without interest; for 2 $\times$ 1{,}117\,MWe net (design data of the units) that is 13{,}876\,USD/kW nominal; with uniform spending over 2013--2023 the CPI-U factor is 1.262, giving 17{,}516\,USD/kW in 2025 dollars and 15{,}656\,EUR/kW, 0.7\% above the 15{,}540 used. The extremes, 30 and 35 billion, give 15{,}150 and 17{,}680\,EUR/kW. With 15 years of construction the risk case has fixN 878.3\,EUR/kW per year at 5\% and 1{,}334.2 at 7\%; its results were not computed but read off the curve by linear interpolation between computed points. Construction takes 7--15 years, against an IAEA median of 5.7--6.3 years \citep{Ritchie2023}. Hinkley Point C, estimated at about 35 billion pounds of 2015 for 3.26\,GW, is not in the median.

\paragraph{Nuclear: other parameters.} Life is 80 years, as already licensed for 25 US units \citep{NRCslr}. Operation and fuel are 27\,EUR/MWh at 90\% load, consistent with fixed O\&M of 136--158\,USD/kW-yr \citep{Lazard2025} plus fuel, and 5\,EUR/MWh for decommissioning and waste, within the fee of 4.7--8.5\,EUR/MWh estimated in Sweden \citep{Riksgalden2025}. Stand-alone, the cost of a nuclear MWh is $\mathrm{lcoe}=K\phi\,\mathrm{CRF}(r,80)/(0.9\times8760)\times1000+32$. SMR cases use a 5-year build and 60-year life with the contract costs of Darlington (17{,}200\,EUR/kW for unit 1, 11{,}700 for four units \citep{POWERDarlington2025}) and two values that are not contracts: 6{,}250\,EUR/kW, the midpoint of the 5{,}000--7{,}500 range given for Rolls-Royce \citep{WNNWylfa2025}, and 3{,}500, a press estimate \citep{Sole24SMR}.

\paragraph{Long-duration and hydrogen storage.} The 100-hour storage costs 20\,USD/kWh over 20 years at 5\%, 1{,}486\,EUR per MWh of capacity per year with 50\% efficiency \citep{FormTarget2021,FormRTE2026}; life 20 years and O\&M of 20\% are our assumptions. It is the developer's target, not an observed price: no iron-air plant is in commercial operation and no current cost is disclosed \citep{FormMIT2026}. Hydrogen is modelled with electrolysis, geological storage and turbines. Large-scale geological hydrogen storage exists only as studies and pilots, so its costs are first-plant (PI: electrolyser 1{,}819\,EUR/kW, storage 113\,EUR/MWh-yr, turbine 1{,}079\,EUR/kW) or target (OB: 450, 56 and 1{,}324) values from studies \citep{HystoriesD86,BaringaDESNZ2025,IEAGHR2025}, not observed prices. Lithium-ion batteries are priced at the 2025 MACSE auction, an observed contract price. Imports are priced at 115.94\,EUR/MWh, the 2025 average PUN \citep{GME2025}, under the rule described in Section~\ref{sec:model}; in the variant of the earlier version the price is raised by 17\% (our assumption) in the three critical weeks of each year, with availability and credit of 38\%. The gas price in the NZ variants is 38.5\,EUR/MWh, the 2025 IG Index of the same report \citep{GME2025}.

\begin{table}[htbp]
\centering\footnotesize
\caption{Model inputs and their status. DATA: read in a source; DERIVED: calculation on data; ESTIMATE: author's judgement with a declared method; ASSUMPTION: model choice without a source, tested in a sensitivity.}\label{tab:inputs}
\input{tab_inputs}
\end{table}

\subsection{Time step, weather years and checks}\label{sec:stepmethods}
The six-hour step underestimates the value of nuclear (Section~\ref{sec:step}). Central values are therefore three-hour results. For the PI case we add a Richardson extrapolation of the gap towards a one-hour step. We assume $\mathrm{gap}(h)=G_0-c\,h^p$ with $h$ in hours; from the 12-, 6- and 3-hour gaps, $2^p$ is the ratio of the increments 12 to 6 and 6 to 3 hours, $c\,3^p=(\text{increment }6\to3)/(2^p-1)$, $G_0=\mathrm{gap}(3)+c\,3^p$ and $\mathrm{gap}(1)=G_0-c$. The extrapolation is accepted when the ratio of the increments lies between 0.35 and 0.65.
To test years outside the base sample, we rank the 38 PECD years 1982--2019 by the mean of the three worst weeks of the combined index of 171\,GW of solar and 44.6\,GW of onshore wind. The four worst years are 2015, 1998, 2014 and 1994 (13.58 to 14.66\,GW against 16.3 to 19.1\,GW for 2016--2019). They replace the weather of 2016--2019, paired in that order with the load of those years, at six hours. The criterion ignores cold spells and hydropower, and old weather is paired with recent demand structure.
\FloatBarrier
